%==============================================================================
% Appendix: Proof of Theorem 4 — Shock-Location Admissibility
% Auto-converted from Docs/proof/Theorem 4 draft.md
%==============================================================================
\begin{theorem}[Shock-location admissibility]
\label{thm:shock-location}
In the limit $\sigtau \to 0$ (i.e., the final step of inference), the solution produced by the
proposed sampler satisfies the Rankine--Hugoniot condition at each shock location~$x_{s}(t)$
almost everywhere,
\begin{equation}
\label{eq:rh-appx}
\dot x_{s}(t) = \frac{\flux(u_{L}) - \flux(u_{R})}{u_{L} - u_{R}},
\end{equation}
and the Lax entropy condition $\flux'(u_{L}) \ge \dot x_{s} \ge \flux'(u_{R})$.
\end{theorem}

\begin{proof}

\noindent\textbf{Assumptions.}

\begin{enumerate}
\item \textbf{(Perfect score matching at the interfacial layer.)}
  The BV-aware score parameterization~\eqref{eq:bv-aware} is optimized to global optimality
  in a neighborhood of the discontinuity, i.e.\ $\sth(u,\tau) \equiv \nabla_u \log \ptau(u)$.
  We thus analyze the true score field.

\item \textbf{(Single-shock target manifold.)}
  The target distribution $\rhotphys(u)$ concentrates on a submanifold $\mathcal{M}_{t}$
  in function space (or its discretization~$\R^{N}$).%
  For the sampling trajectory under consideration, the corresponding target
  Kruzhkov entropy solution $\physsol(x,t)$ at the current time~$t$ possesses a single
  isolated shock located at~$x_{s}(t)$, with left and right limits
  \begin{equation}
  u_{L} \coloneqq \physsol(x_{s}^{-},t), \qquad
  u_{R} \coloneqq \physsol(x_{s}^{+},t).
  \end{equation}

\item \textbf{(Local planarity in high dimension.)}
  Since the grid dimension~$N$ is large, in the limit $\tau\to0$ the separating
  manifold of data modes (i.e., the locus where the score exhibits a shock) can be
  locally approximated as a hyperplane.
  Its signed distance function~$\phissh(u,\tau)$ satisfies
  $\abs{\nabla_u \phissh} = 1$.

\item \textbf{(Compressive score gradient.)}
  When the reverse sampling trajectory $u(\tau)$ approaches the target
  manifold~$\mathcal{M}_{t}$, the local flow induced by the network
  parameterization~$\sth$ is compressive along the shock
  normal~$\nabla\phissh$:
  \begin{equation}
  \partial_{\phissh} \sth \cdot \nabla\phissh < 0.
  \end{equation}
  This guarantees that the score field acts as a centripetal force, pushing
  probability mass toward the interfacial manifold.
\end{enumerate}

%==========================================================================
% Part I: Geometric Correspondence via the Viscosity-Matched Schedule
%==========================================================================
\medskip
\noindent\textbf{Part I: Geometric correspondence via the viscosity-matched schedule.}

Let the forward noising process be a variance-exploding SDE.
Under the schedule~\eqref{eq:visc-matched}, the noise variance is tied to the physical
viscosity~$\physvis$:
\begin{equation}
\sigtau^{2} = 2\physvis\,\tau, \qquad \tau \in [0,T_{d}].
\end{equation}
The forward SDE then reads
\begin{equation}
du = \sqrt{2\physvis}\; dw.
\end{equation}
The corresponding Fokker--Planck equation for the density $p(u,\tau)$ in data space~$\R^{N}$ is
\begin{equation}
\label{eq:fp}
\dtau p(u,\tau) = \physvis\,\Lap_{u} p(u,\tau).
\end{equation}

Differentiating the score $\score(u,\tau) = \nabla_u \log p = \nabla_u p / p$ with
respect to~$\tau$ and invoking~\eqref{eq:fp} gives the exact evolution of the score:
\begin{align}
\dtau \score
  &= \physvis\,\nabla_u\!\Bigl( \frac{\Lap_u p}{p} \Bigr) \notag\\
  &= \physvis\,\Lap_u \score + \physvis\,\nabla_u(\abs{\score}^{2}).
\end{align}
Since $\score$ is a gradient field, $\nabla_u(\abs{\score}^{2}) = 2(\score\cdot\nabla_u)\score$.
Note that the relation $\score = \nabla_u\log p$ is essentially the classical
\emph{inverse Cole--Hopf transform}, which explicitly converts the linear heat equation
(Fokker--Planck) into a nonlinear PDE.
We thus obtain the \textbf{high-dimensional viscous Burgers equation at the score level}:
\begin{equation}
\label{eq:score-burg}
\dtau \score + (-2\physvis\,\score\cdot\nabla_u)\,\score = \physvis\,\Lap_u \score.
\end{equation}

\begin{lemma}[Effective viscosity equivalence]
\label{lem:visc-equiv}
Via the schedule~\eqref{eq:visc-matched}, the effective viscosity of the score dynamics is
numerically identical to the physical viscosity~$\physvis$ of the target equation.
\end{lemma}

The target PDE with small physical viscosity is
\begin{equation}
\label{eq:target-pde}
\dt u(x,t) + \dx \flux(u(x,t)) = \physvis\,\partial_{xx} u(x,t).
\end{equation}
Under discretization, the components of $u \in \R^{N}$ satisfy $u_i \approx u(x_i,t)$.
The target distribution $\rhotphys(u)$ (the boundary condition at $\tau=0$) is the physical
weak-solution distribution.
In the limit $\tau\to0$, the reverse sampling trajectory $u(\tau)$ is pulled by the
score~$\score(u,\tau)$ toward the support manifold~$\mathcal{M}_{t}$ of~$\rhotphys(u)$.

The BV-aware score parameterization~\eqref{eq:bv-aware} reads
\begin{equation}
\sth(u,\tau)
  = \nabla\phisbg + \frac{\jumpamp}{2}\tanh\!\Bigl(\frac{\phissh}{2}\Bigr)\nabla\phissh.
\end{equation}
Near the interface, the smooth background $\nabla\phisbg$ is negligible, and the true
score exhibits the asymptotic inner-scale form
\begin{equation}
\label{eq:score-inner}
\score(u,\tau)
  \sim \frac{\jumpamp(u)}{2}
       \tanh\!\Bigl( \frac{\operatorname{dist}(u,\mathcal{M}_{t})}
                          {\sqrt{2\physvis\,\tau}} \Bigr)
       \mathbf{n}(u),
\end{equation}
where $\mathbf{n}(u) = \nabla_u \operatorname{dist}(u,\mathcal{M}_{t})$ is the normal vector
to the manifold.

\begin{lemma}[Geometric correspondence in the double limit]
\label{lem:geom-corresp}
When $\tau\to0$ and $\physvis\to0$ hold simultaneously, the $\tanh$ profile
in~\eqref{eq:score-inner} degenerates to a step function.
Consequently, the location where the score field exhibits a shock is determined in data space
by $\operatorname{dist}(u,\mathcal{M}_{t}) = 0$.
Mapped back to physical space, this manifold boundary coincides exactly with the grid
coordinates where $u(x,t)$ exhibits a jump, i.e., the physical shock position~$x_{s}(t)$.
\end{lemma}

%==========================================================================
% Part II: Weak-Form Limit and Rankine--Hugoniot Condition
%==========================================================================
\medskip
\noindent\textbf{Part II: Weak-form limit and derivation of the Rankine--Hugoniot condition.}

\smallskip
\noindent\textit{2.1  Tweedie profile near the shock.}

For the reverse diffusion, the clean physical state $\uhat_{0}(x,t)$ is estimated via the
Tweedie formula
\begin{equation}
\uhat_{0}(x,t) = u(\tau) + \sigtau^{2}\,\score(u,\tau).
\end{equation}
Under the schedule~\eqref{eq:visc-matched} and the BV-aware parameterization, restricted
to a one-dimensional physical section, the signed distance expands as
$\displaystyle \phissh \approx \frac{x - x_{s}(t)}{\sqrt{2\physvis\,\tau}}$.

\begin{lemma}[Viscous shock profile]
\label{lem:visc-profile}
In a neighborhood of $x_{s}(t)$, the Tweedie-reconstructed physical field
$\uhat^{\tau}(x,t)$ possesses a smooth, monotone shock transition layer.
To make the physical scales explicit, define the layer thickness
$\varepsilon = \sigtau = \sqrt{2\physvis\,\tau}$.
It is described by
\begin{equation}
\label{eq:visc-profile}
\uhat^{\tau}(x,t)
  \approx \frac{u_{L}+u_{R}}{2}
        - \frac{u_{L}-u_{R}}{2}
          \tanh\!\Bigl( \frac{x - x_{s}(t)}{2\varepsilon} \Bigr).
\end{equation}
Thus $\uhat^{\tau}$ is strictly between $u_{L}$ and $u_{R}$ and has controlled total variation.
\end{lemma}

\smallskip
\noindent\textit{2.2  Godunov guidance and weak form.}

During generation, an entropy-compatible Godunov flux difference is used as PDE guidance.
Thus, as $\tau\to0$, the generated sequence $\uhat^{\tau}$ approximately satisfies a weak form:
for every smooth test function $\varphi \in C_{c}^{\infty}(\R\times(0,T))$,
\begin{equation}
\label{eq:weak-form}
I(\varepsilon)
  = \iint_{\R\times\R^{+}}
    \bigl( \uhat^{\tau}\,\partial_{t}\varphi + \flux(\uhat^{\tau})\,\partial_{x}\varphi \bigr)
    \,dx\,dt
  = \mathcal{O}(\varepsilon).
\end{equation}

Because the $\tanh$ profile guarantees $\uhat^{\tau} \in L^{\infty} \cap \BV$, the dominated
convergence theorem allows the limit $\varepsilon\to0$ to be taken inside the integral.
\textbf{This step is particularly emphasized: the local monotonicity induced by
parameterization~\eqref{eq:bv-aware} is the key topological constraint that renders this weak
limit legitimate.}
Traditional diffusion models (e.g., those forcing a Gaussian fit across shock interfaces)
inevitably produce Gibbs oscillations, causing the~$\BV$ norm to blow up in the limit and
precluding application of dominated convergence.
The profile~\eqref{eq:visc-profile} converges strongly to the piecewise-constant function
\begin{equation}
\physsol(x,t) =
\begin{cases}
u_{L}, & x < x_{s}(t),\\[2pt]
u_{R}, & x > x_{s}(t).
\end{cases}
\end{equation}

Partition the spacetime domain along the shock trajectory $\Gamma\!: x = x_{s}(t)$ into a
left region $D^{-}$ ($x < x_{s}$) and a right region $D^{+}$ ($x > x_{s}$).
Then
\begin{equation}
\label{eq:limit-split}
0 = \lim_{\varepsilon\to0} I(\varepsilon)
  = \iint_{D^{-}} \bigl( u_{L}\,\partial_{t}\varphi + \flux(u_{L})\,\partial_{x}\varphi \bigr)
    \,dx\,dt
  + \iint_{D^{+}} \bigl( u_{R}\,\partial_{t}\varphi + \flux(u_{R})\,\partial_{x}\varphi \bigr)
    \,dx\,dt.
\end{equation}
Inside each $D^{\pm}$, the solution $u_{L},u_{R}$ is a strong solution of the conservation law.
Applying Green's theorem to transfer derivatives from $\varphi$ onto~$u$,
\begin{equation}
\label{eq:green}
\begin{aligned}
\iint_D \bigl( u\,\partial_{t}\varphi + \flux(u)\,\partial_{x}\varphi \bigr) \,dx\,dt
&= \int_{\partial D}
   \varphi\,\bigl( u\,n_{t} + \flux(u)\,n_{x} \bigr)\,ds \\
&\quad - \iint_D \varphi\,\bigl( \partial_{t}u + \partial_{x}\flux(u) \bigr) \,dx\,dt.
\end{aligned}
\end{equation}
Since the PDE holds pointwise in $D^{\pm}$, the area term vanishes.
Because $\varphi$ has compact support, only the boundaries along~$\Gamma$ contribute.
On $\Gamma$, the outward normal (pointing toward $D^{+}$) is
\begin{equation}
\mathbf{n} = (n_{x}, n_{t})
           = \frac{1}{\sqrt{1 + \dot x_{s}^{2}}}
             \bigl( 1,\, -\dot x_{s} \bigr).
\end{equation}
Accounting for orientation, the jump condition emerges:
\begin{equation}
\label{eq:jump-cond}
\int_{\Gamma} \varphi \cdot
\Bigl[
  \bigl( u_{L}(-\dot x_{s}) + \flux(u_{L}) \bigr)
  - \bigl( u_{R}(-\dot x_{s}) + \flux(u_{R}) \bigr)
\Bigr]
\frac{1}{\sqrt{1 + \dot x_{s}^{2}}}
\,ds = 0.
\end{equation}
Since $\varphi$ is arbitrary along $\Gamma$, the algebraic bracket must vanish identically,
yielding
\begin{equation}
\dot x_{s}\,(u_{L} - u_{R}) = \flux(u_{L}) - \flux(u_{R}).
\end{equation}
For a genuine shock $u_{L} \neq u_{R}$, division by the difference gives the
Rankine--Hugoniot condition
\begin{equation}
\label{eq:rh-derived}
\dot x_{s}(t) = \frac{\flux(u_{L}) - \flux(u_{R})}{u_{L} - u_{R}}.
\qquad\square
\end{equation}

%==========================================================================
% Part III: Entropy Compatibility and the Lax Condition
%==========================================================================
\medskip
\noindent\textbf{Part III: Entropy compatibility and the Lax condition.}

\smallskip
\noindent\textit{Background: the non-uniqueness of weak solutions.}
The Rankine--Hugoniot condition derived above is a \emph{necessary} but not sufficient
condition for a weak solution.
Physically inadmissible ``expansion shocks'', which dissipate energy in an
entropy-violating manner, also satisfy the R--H condition in nonlinear evolution models.
Only by imposing a rigorous entropy condition---the Lax entropy condition, motivated by
physical mechanisms such as the Double-Burgers coupling---do we ensure that the generated
solution is not merely a ``numerical fit'' to data, but a genuine approximation of the
\emph{physically correct} conservation law.

\smallskip
\noindent\textit{3.1  Divergence--characteristic relation.}

For a convex flux $\flux$ (e.g., the Burgers flux $\flux(u)=u^{2}/2$), a physical shock
must satisfy the Lax entropy condition
\begin{equation}
\label{eq:lax}
\flux'(u_{L}) > \dot x_{s} > \flux'(u_{R}).
\end{equation}
In the reverse-time generative ODE,
\begin{equation}
\frac{du}{d\tau} = -\frac12\,\beta(\tau)\bigl[ u + \sth(u,\tau) \bigr],
\end{equation}
Tweedie reconstruction yields
$\partial_{x}\uhat \approx \partial_{x}(u + \sigtau^{2}\,\sth) \to -\infty$
near $x_{s}$.
This singular negative gradient corresponds to the compression of characteristics
toward the shock in physical space.

\smallskip
\noindent\textit{3.2  Kruzhkov entropy inequality.}

To prove that $\uhat$ is an entropy solution, we verify the Kruzhkov inequality for every
constant~$k$ (in the sense of distributions):
\begin{equation}
\label{eq:kruz}
\partial_{t}\,\abs{\uhat - k}
  + \partial_{x}\bigl[ \operatorname{sgn}(\uhat - k)\,(\flux(\uhat) - \flux(k)) \bigr]
  \le 0.
\end{equation}
By Theorem~1 (Double-Burgers coupling), the score field itself satisfies a viscous
Burgers equation.
Hence the generative dynamics driven by $\sth$ constitute a JKO gradient flow that
minimizes an energy functional.
A jump $u_{L} \to u_{R}$ violating the Lax condition (i.e., $u_{L} < u_{R}$) would
correspond to an expansion (rarefaction) fan whose characteristics diverge outward.
Under the BV-aware parameterization~\eqref{eq:bv-aware}, the sign of the score gradient is
determined by the $\tanh$ profile.

\smallskip
\noindent\textit{3.3  Proof by contradiction.}

Suppose a non-physical shock (an expansion shock) satisfies the R--H condition but has
$u_{L} < u_{R}$.

\begin{enumerate}
\item[\textbf{(i)}] \textbf{Score behavior.}
  The Tweedie formula then requires the score~$\sth$ to supply an outward ``push'' in order
  to sustain this unstable discontinuity.

\item[\textbf{(ii)}] \textbf{Burgers consistency loss.}
  The loss term $\Lburg$ (Burgers consistency) enforces
  $\dtau \score + (\score\cdot\nabla)\score = \physvis\,\Lap\score$.
  An expansion jump would be immediately regularized by the diffusion term~$\Lap\score$,
  producing a large residual in~$\Lburg$.

\item[\textbf{(iii)}] \textbf{Conclusion.}
  Minimizing the total loss systematically eliminates all non-physical jumps, because only
  compressive configurations satisfying the Lax condition~\eqref{eq:lax} are compatible with
  the stable evolution operator of~$\Lburg$.
\end{enumerate}

\smallskip
\noindent\textit{3.4  Stability in the $\tau\to0$ limit.}

Finally, the Oleinik entropy condition provides the estimate near the shock:
\begin{equation}
\frac{\uhat(x+a,t) - \uhat(x,t)}{a} \le \frac{C}{t},
\end{equation}
which is guaranteed by the controlled width of the $\tanh$ layer in
parameterization~\eqref{eq:bv-aware}.
As $\tau$ decreases, the profile steepens but remains strictly monotone whenever $u_{L}>u_{R}$,
thereby precluding non-physical oscillations.

\medskip
\noindent
Combining the three parts above, the generated solution satisfies both the Rankine--Hugoniot
condition and the Lax entropy condition almost everywhere.
\end{proof}
